% Faithful conversion of ages-i-sprint-2; original XML and source characters retained.
\hypertarget{__RefHeading___Toc4747_3118176633}{}
\refstepcounter{subsection}
\subsection*{2.3.3 Sprint 2\originalspace{}}
\label{unit:ages-i-sprint-2}
\addcontentsline{toc}{subsection}{2.3.3 Sprint 2\originalspace{}}

% text-p-0509
\noindent \originaltab{}Para a Sprint 2, desenvolvi a última tela necessária para o Tablet. Fui convidado por meu colega de equipe a participar da criação da tela de suplementos escolhidos. Esta, tem uma grande importância para o projeto, pois nela o usuário visualizará quais são os suplementos gerados pelo nosso algoritmo. Devido a isso, me empenhei ao máximo para criar uma tela harmoniosa. A realização dessa task se deu de forma árdua, visto que muitas alterações tiveram de ser feitas (será explicado na caixa de comentários).\originalspace{}

% text-p-0510
\noindent \originaltab{}Além disso, decidi me “aventurar” e explorar tasks relativas à API do projeto. Encontrei uma tarefa chamada “criar endpoints para adicionar e remover suplementos do estoque de lojista”, e a designei a mim, sem mesmo saber o que eram endpoints. Inicialmente, precisei de auxílio do meu colega de equipe Matheus Chong para configurar o Banco de Dados no Docker e rodar o Back-end corretamente. Após isso, estudei sobre conceitos e como aplicar os enpoints, além de consultar meus colegas AGES III para a realização da task. Também compreendi melhor o funcionamento do nosso Banco de Dados através do Prisma Studio. Adquiri o conhecimento necessário e pude desenvolver um código consiso e eficiente.

% text-p-0511
\noindent \originaltab{}Tive muitos aprendizados durante a Sprint 2. Embarquei em novos desafios, consultei colegas experientes para me tornar mais competente, explorei tecnologias novas, e apliquei diversos conhecimentos adquiridos nas aulas do curso de Engenharia de Software. Vale destacar a cadeira de Gerência de Configuração de Software, que se demonstrou muito útil na prática de gerenciamento e versionamento de código. Pude aprender conceitos de Git ainda mais complexos durante a Sprint, o que só foi possível graças a sólida base fornecida por essa cadeira.

\begin{reflection}
% text-p-0521
\noindent \originaltab{}O desenvolvimento da tela de suplementos escolhidos foi um pouco estressante. Meu colega de equipe havia me chamado para ajudá-lo com a task já na metade de seu desenvolvimento. Quando visualizei o código-fonte, percebi que ele havia ignorado o modo de construção de telas modulares. Ou seja, todos os componentes que foram desenvolvidos pela equipe durante a última Sprint, visando facilitar e agilizar o trabalho de construção das telas, não estavam sendo utilizados.

% text-p-0522
\noindent \originaltab{}Devido a isso, tivemos de reimplementar os componentes modulares Header, Scrollwrapper, além de usar as cores padronizadas, evitando redundâncias nos arquivos. Expliquei a ele que deveríamos seguir uma série de regras para otimizar o desenvolvimento e facilitar o compartilhamento de código entre a equipe. Eu também havia cometido os mesmo erros na Sprint anterior, e para evitar que meu colega tivesse de refazer sua tarefa do zero, incentivei-o a analisar os componentes que havíamos criado anteriormente para que ele pudesse aproveitá-los no desenvolvimento de suas próximas tasks.

% text-p-0523
\noindent \originaltab{}Outro ponto que gostaria de comentar nesta caixa de texto foi a minha iniciativa de me designar uma task relativa à API do projeto. Estava um pouco apreensivo, pois estaria envolvido com outro repositório, que possuía arquivos diferentes do que eu havia trabalhando anteriormente, e com tasks que, até o momento, estavam fora da minha compreensão. Apesar desses percalços, estudei e consultei meus colegas. E, rapidamente, me tornei apto para realizar as tarefas do Backend.\originalspace{}\originalspace{}

\end{reflection}
